\documentclass[11pt]{article}
\usepackage[a4paper,margin=25mm]{geometry}
\usepackage{amsmath,amssymb,amsthm,mathtools,booktabs,hyperref,xurl}
\setlength{\emergencystretch}{2em}
\newtheorem{theorem}{Theorem}
\newtheorem{lemma}[theorem]{Lemma}
\newtheorem{proposition}[theorem]{Proposition}
\newtheorem{corollary}[theorem]{Corollary}
\newtheorem{remark}[theorem]{Remark}
\newcommand{\Z}{\mathbb Z}\newcommand{\Q}{\mathbb Q}
\title{A corrected quartic source from Jones's binomial theorem}
\author{\(\mathbb Z\mathbb Q\) Diophantine Project}
\date{23 November 2026}
\begin{document}
\maketitle

\begin{abstract}
The previous quartic tables were compiled from an emended display of
Jones's 1980 Theorem~3. Their precise source universality was left as
the unverified premise \(\mathcal U_{\rm em}\). Here we replace that
source, rather than asserting that the old display was correct.
Jones's preceding Theorem~2 is an exponential--binomial coding theorem.
The independently proved binomial certificate supplies its polynomial
replacement, with the required exponent \(n^6\), and gives an exact
equivalence to that classical theorem. Only two quadratic residuals
change in each of the existing tables; neither the coordinate count nor
the number of residuals increases. The new source uses Jones's
Theorem~2 as an explicitly identified external result. We do not claim
to have independently reproved its coding argument, or the arithmetic
premises of the project's rational integer predicates.
\end{abstract}

\section{The precise external theorem}

Fix positive integers \(x,z,u,y\), and set
\[
 L=5^{60},\qquad B=b^5,\qquad C=1+xB+g.
\]
To avoid confusing the source's unrelated uses of a small letter, write
\(\omega\) for an exponent and \(\nu\) for a binomial quotient.
Jones's 1980 Theorem~2 \cite[p.~860]{Jones80} asserts that
\(x\in W_{z,u,y}\) if and only if the following relations have positive
integer witnesses
\[
 b,e,g,\ell,m,n,q,r,t,\omega,\alpha,\nu,\theta,\lambda:
\]
\begin{align}
 e\ell g^2+\alpha&=(b-xy)q^2,&q&=b^L,\label{prefix1}\\
 \lambda+q^4&=1+\lambda B,&\theta+2z&=B,\label{prefix2}\\
 \ell&=u+t\theta,&e&=y+m\theta,\label{prefix3}\\
 b&=2^\omega,&n&=q^{16},\label{exp}\\
 r&=S(n^2-n)+T(n^2-1),&\nu n^2&=\binom{2r}{r},\label{binom}
\end{align}
where the two displayed polynomial abbreviations are
\begin{align}
 S&=g+eq^3+\ell q^5+
       \bigl(2(e-z\lambda)C^4+\lambda B+\lambda Bq^4\bigr)q^4,
       \label{S}\\
 T&=q^3-b\ell+\ell+\theta\lambda q^3+(B-2)q^5.
       \label{T}
\end{align}
The family \(W_{z,u,y}\) enumerates the computably enumerable sets.
In particular, some fixed \(z,u,y\) give a nonrecursive set.

The source being cited is \emph{Theorem~2}, not the faulty expanded
Theorem~3. The 1980 note states its theorems without proofs; the
paper \cite{Jones82} is the later full treatment. Our use of the
classical coding theorem is an external input, not a claim of a new
independent proof. A public transcription of the original four-page
note was checked at \cite{Transcription}; it reproduces the coding
prefix and its single central-binomial condition. The exact polynomial
replacement below is proved from the project's separate, complete Pell
certificate, independently of any source emendation.

\section{A corrected polynomial replacement}

For an integer \(a\ge2\), define Pell sequences by
\[
 X_j(a)+Y_j(a)\sqrt{a^2-1}
       =(a+\sqrt{a^2-1})^j,\qquad j\ge0.
\]
The following local result has a complete elementary proof in
\cite[Theorem~1]{Binomial}.

\begin{lemma}[The independently proved local certificate]\label{local}
Let \(r,M,b\in\Z\) with \(r\ge1\), \(M\ge2\), and
\(2\le b\le2M\). Then
\[
 b\text{ is a positive power of }2,
 \qquad M\mid\binom{2r}{r}
\]
is equivalent to the existence of positive integer witnesses
\(w,s,a,p,k,\tau,h,c,d,\varphi,\gamma,\eta,i,F,j,o\)
satisfying
\begin{align}
 a&=(wM+1)rsM,&p&=2ws^2r^2M^3,\label{ap}\\
 k&=r+1+h(p-1),&\tau^2&=(p^2-1)k^2+1,\label{small}\\
 4(c-ksM)^2+\eta&=k^2,&c&=2r+1+\varphi,\label{close}\\
 d&=bw+(a-2)c+(4a-5)\gamma,\label{linear}\\
 d^2&=(a^2-1)c^2+1,\label{norm1}\\
 F^2&=(a^2-1)i^2c^4+1,\label{norm2}\\
 (d+oF)^2&=(G^2-1)(2r+1+jc)^2+1,\label{norm3}\\
 G&=a+F^2(c+1-a).\label{bridge}
\end{align}
The alternative bridge \(G=a+F^2(d^2-a)\) is equivalent as
well. Every positive solution forces
\[
 bw=2^{2r+1},\qquad c=Y_{2r+1}(a),\qquad d=X_{2r+1}(a),
 \qquad k=Y_{r+1}(p).
\]
\end{lemma}

\begin{theorem}[A corrected classical universal source]\label{corrected}
Retain \eqref{prefix1}--\eqref{prefix3}, the equation \(n=q^{16}\),
and the equation for \(r\) in \eqref{binom}, with \(S,T\) defined
by \eqref{S}--\eqref{T}. Delete \(\omega,\nu\), the exponential
equation \(b=2^\omega\), and the binomial equation in
\eqref{binom}. Adjoin the polynomial equations in
Lemma~\ref{local} with \(M=n^2\).

The resulting system has 28 positive integer unknowns and is solvable
if and only if \(x\in W_{z,u,y}\). Its Pell parameter is
\begin{equation}\label{n6}
                         p=2ws^2r^2n^6.
\end{equation}
This theorem uses Jones's Theorem~2 and Lemma~\ref{local}; it does
not use \(\mathcal U_{\rm em}\), \(\mathsf U_{MT}\), or a
literal transcription of either expanded historical source.
\end{theorem}
\begin{proof}
First suppose that the relations of Jones's Theorem~2 hold. The equation
\(b=2^\omega\) gives \(b\ge2\). Because \(q=b^L\) and
\(n=q^{16}\), we have \(n\ge b\), whence
\[
 n^2\ge2,\qquad 2\le b\le2n^2.
\]
The positive binomial quotient \(\nu\) gives
\(n^2\mid\binom{2r}{r}\). Lemma~\ref{local} now supplies all
sixteen new witnesses, and every unchanged equation remains true.

Conversely, suppose that the corrected polynomial system has positive
integer witnesses. The positivity of \(\theta=b^5-2z\), with
\(z\ge1\), implies \(b\ge2\). Again \(n\ge b\), so all
hypotheses of Lemma~\ref{local} hold with \(M=n^2\). The lemma
gives an integer \(\omega\ge1\) such that \(b=2^\omega\), and
the quotient
\[
                        \nu=\binom{2r}{r}/n^2
\]
is a positive integer. These restore exactly the two deleted relations
of Theorem~2. Therefore \(x\in W_{z,u,y}\).

The original theorem has fourteen unknowns; we delete two and adjoin
sixteen, giving \(14-2+16=28\). The symbols \(G,S,T,B,C,M\)
in this statement are polynomial abbreviations, not additional
unknowns. Equation \eqref{n6} is \eqref{ap} with \(M=n^2\).
\end{proof}

\begin{remark}[What this does and does not repair]
The theorem establishes a new, corrected complete source. It does not
assert that replacing \(n^2\) by \(n^6\) in the old system preserves
that old system's solution set. In particular, it does not prove the
repository's old premise \(\mathcal U_{\rm em}\). The historical
tables retain their previous statements and their previous qualification.
\end{remark}

\section{Two residual changes in the complete quadratic tables}

The current tables use an already stored coordinate \(n_2=n^2\).
Their six-coordinate local block, after the earlier elimination of
a temporary product, is
\begin{align*}
 Y&=sn_2,&L_0&=sr,&L_{0,2}&=L_0^2,&U&=wn_2,\\
 p&=2UL_{0,2},&&2L_0a=(p+2L_{0,2})n_2.
\end{align*}
Here the literal register names in the coefficient tables are
\(\mathtt{sr}=L_0\), \(\mathtt{sr2}=L_{0,2}\).

\begin{proposition}[A division-free corrected block]\label{two}
Reuse the same six coordinates \(Y,L,L_2,U,p,a\), but replace the
two rows defining \(L_0\) and \(a\) by
\[
                         L=rY,\qquad a=(U+1)L.
\]
Retain the other four rows, interpreting the reused registers as
\(\mathtt{sr}=L\), \(\mathtt{sr2}=L_2\). The resulting block is
\begin{equation}\label{block}
 Y=sn_2,\quad L=rY,\quad L_2=L^2,\quad U=wn_2,
 \quad p=2UL_2,\quad a=(U+1)L.
\end{equation}
It implies exactly
\[
                  a=(wn^2+1)rsn^2,
                  \qquad p=2ws^2r^2n^6.
\]
Every row has total degree at most two. No coordinate, row, integer
predicate, reciprocal guard, or positivity witness is added.
\end{proposition}
\begin{proof}
Direct substitution in \eqref{block} gives
\(L=rsn^2\), \(L_2=r^2s^2n^4\), and both required formulas.
Conversely, the defining values of the six coordinates satisfy every
row. No division is made. An exact occurrence audit of each complete
predecessor table verifies that the register \(\mathtt{sr}\)
occurs only in its defining row, its square row, and the row defining
\(a\). Its square occurs in its definition, the row for \(p\),
and the old row for \(a\). Thus the changes above have no unaccounted
effect elsewhere in the complete system.
\end{proof}

The existing twelve integrality predicates, the cofinite instance
restriction \(x>2z\), and the reconstruction of every other integer
coordinate are unchanged. In particular the new register \(L=rY\)
is an integer when its predecessor inputs are integers. The signed-Pell
and positivity reductions of \cite{Pell} explicitly cover both the
\(n^2\) and \(n^6\) choices of \(p\). They therefore apply to
the corrected block. The row
\[
                        \zeta^2=(a-1)(d-1)
\]
continues to choose the positive root, and the new bridge
\eqref{bridge} is legitimate. All coefficient-sharing and rational
rechart transformations outside this local block are preserved
literally.

The reciprocal guard has a noncircular reconstruction order. At the
stage when the positive integer inputs \(s,n\) have been recovered,
\(Y=sn^2\ge4\). The retained equation
\(\eta=k^2-4(c-kY)^2>0\) first gives \(k\ne0\). If \(c=0\),
it would instead give
\(\eta=(1-4Y^2)k^2\le0\), a contradiction. Thus \(c\ne0\)
without using the later Pell conclusion \(c>2r+1\). The row
\(vE=c^2\) now forces \(v\ne0\), so the reciprocal integer
predicate legitimately recovers \(i=1/v\in\Z\).
The monic norm equations then give the required integrality of
\(F\) and \(\tau\), before the signed-Pell argument reconstructs
their positive representatives and the remaining positive source
witnesses.

\begin{corollary}[Corrected quartic transfers]\label{transfer}
Use any of the project's current twelve-input integrality interfaces,
with all of its stated arithmetic hypotheses. Apply
Proposition~\ref{two} to the complete tables of v92. For the fixed
coefficients \(z,u,y\) supplied by Jones's Theorem~2, the resulting
family has undecidable rational solvability. The counts are
\begin{center}
\begin{tabular}{lrr}
\toprule
Interface & Rational unknowns & Quadratic residuals\\
\midrule
General &1103&953\\
Specified &1283&1109\\
Ordinary rational \(abc\)&837&725\\
\bottomrule
\end{tabular}
\end{center}
The instance \(x\) occurs in the coefficients. Squaring and summing
the residuals produces a polynomial of degree exactly four.
The source premise \(\mathcal U_{\rm em}\) is absent for these
new tables; Jones's stated coding theorem is the external source.
\end{corollary}
\begin{proof}
Rational soundness of the integer interfaces and the existing
reconstruction prove that every rational zero gives a positive
integer solution of the corrected source, after reconstructing the
signed Pell witnesses as in \cite{Pell}. Theorem~\ref{corrected}
therefore gives membership in the chosen nonrecursive set.
Conversely, each member has a corrected positive source witness, and
the retained completeness theorems for the shared integer predicates
give a rational zero of the complete table.

Deleting the finitely many instance values \(x\le2z\) does not
make a nonrecursive set recursive: a decision procedure for its tail,
together with a fixed finite lookup table, would decide the original
set. The effective family on this cofinite tail is therefore
undecidable. Proposition~\ref{two} changes no counts. A sum of
rational squares vanishes precisely when every residual vanishes.
At least one residual has a nonzero quadratic homogeneous part;
the sum of the squares of such parts cannot vanish identically over
the reals, so the total degree is exactly four.
\end{proof}

This corollary does not independently establish the project's integer
predicates. In particular ``general'' does not mean that every retained
elliptic, prime-pattern, pole, and height input has been independently
proved here. The small stored source members remain fixtures where the
original interface has an ineffective sufficiently large parameter.
The ceiling-binary power circuit and the ordinary-\(abc\) fixed-offset
contact chart compose with this two-row repair. Their changes do not
touch the two replaced rows. The composed, source-corrected tables have
respectively \(1074\), \(1254\), and \(786\) coordinates, and
\(924\), \(1080\), and \(685\) residuals. The intermediate
ordinary-\(abc\) ceiling table has \(808\) coordinates and \(696\)
residuals. Their independent arithmetic-equivalence proofs remain
separate inputs; the source repair itself saves no coordinates.

\section{Reproducibility and provenance}

Run from the repository root:
\begin{verbatim}
python undecidability/source_repair_v93/build_corrected_source.py \
  --include-compositions
python undecidability/source_repair_v93/verify_corrected_source.py
\end{verbatim}
The builder checks the exact old local rows and the occurrences of
the reused register before modifying the complete sparse coefficient
table. The independent verifier checks the two new polynomial formulas,
compares every unchanged residual by variable name, verifies full
coordinate support and degree, and checks deterministic reproduction
and hashes of all seven complete tables. These are exact algebraic
checks, not a machine proof of Jones's coding theorem or of the
external rational integer predicates.

The mathematical source of the coding prefix is Jones. The local
Pell and binomial proofs and the coordinate-preserving replacement are
project work from v91--v93. The supplied tables preserve the attribution
of every earlier rational integer interface.

\begin{thebibliography}{9}
\bibitem{Jones80} J. P. Jones,
\emph{Undecidable Diophantine equations},
Bulletin of the American Mathematical Society (New Series)
\textbf{3}, no.~2 (1980), 859--862, Theorem~2, p.~860.
\url{https://doi.org/10.1090/S0273-0979-1980-14832-6}.
\bibitem{Jones82} J. P. Jones,
\emph{Universal diophantine equation}, Journal of Symbolic Logic
\textbf{47} (1982), 549--571.
\url{https://doi.org/10.2307/2273588}.
\bibitem{Transcription} Public transcription of the original Jones
1980 note, inspected 23 November 2026:
\url{https://doczz.net/doc/7703447/undecidable-diophantine-equations}.
Used to inspect the primary source's Theorem~2, not as an independent
mathematical authority. The journal citation above identifies the theorem.
\bibitem{Binomial} \(\mathbb Z\mathbb Q\) Diophantine Project,
\emph{An exact Pell certificate for a power of two and central-binomial
divisibility}, 23 November 2026.
\nolinkurl{undecidability/papers/binomial_pell_certificate.tex}.
\bibitem{Pell} \(\mathbb Z\mathbb Q\) Diophantine Project,
\emph{An exact three-Pell index certificate and nine fewer quadratic
coordinates}, 23 November 2026.
\nolinkurl{undecidability/papers/pell_index_optimization.tex}.
\end{thebibliography}
\end{document}
